\section{Map System}
\label{sec:map-system}

\subsection{Overview}

The Map System provides the combat stage for all auto-battles. 
Each Map is an arena-style combat area with a fixed Wave structure and one or more Difficulty levels. 
The MapDevice in the Hideout acts as the entry point: the player selects a Map, a Difficulty (if unlocked) and optionally empowerment modifiers before starting a run.

Maps are tightly integrated with:
\begin{itemize}
	\item the \textbf{Hero \& Skill System} (which team and build enters the Map),
	\item the \textbf{Loot System} (which rewards are generated from enemies and completion),
	\item the \textbf{Research System} (which Maps and Difficulties are available, and which enemy variants exist),
	\item the \textbf{Inventory System} (where all drops are ultimately stored).
\end{itemize}

The Map System itself does not generate items or currencies directly; it defines when and how enemies spawn and which combat events feed into the Loot System.

% RISK: High coupling between Map, Research, Loot and Hero Systems can make iterative balancing slow.
% TODO: Keep MVP Map feature set tight to avoid overdesigning early.

\subsection{Map Types and Layout}

\subsubsection*{Arena Structure}

\begin{itemize}
	\item Maps are \textbf{arena-like}: one main combat area per Map.
	\item There is no free player movement during combat; heroes and enemies move and act autonomously within the arena.
	\item Within this arena, different \textbf{layouts and obstacles} exist per Map:
	\begin{itemize}
		\item terrain features,
		\item obstacles,
		\item lanes and choke points.
	\end{itemize}
	\item Layout influences:
	\begin{itemize}
		\item hero positioning decisions,
		\item pathing and clustering of enemies,
		\item effectiveness of area-based skills and auras.
	\end{itemize}
\end{itemize}

\subsubsection*{Themes and Map Variety}

\begin{itemize}
	\item For MVP, the game targets \textbf{3 high-level themes} (e.g. different fantasy / dungeon styles).
	\item Per theme, there are \textbf{3--5 Maps} with distinct terrain/obstacle layouts:
	\begin{itemize}
		\item visual variation (biomes, architecture),
		\item gameplay variation (different chokepoints, open areas, spawn positions).
	\end{itemize}
	\item In total, MVP aims for \textbf{around 10 Maps} with varied Wave setups and enemy mixes.
\end{itemize}

% RISK: Even arena-only maps with multiple layouts require art and collision work; 10+ maps is ambitious for MVP.

\subsection{Map Selection (MapDevice)}

\subsubsection*{Selection Flow}

From the MapDevice in the Hideout, the player selects:

\begin{itemize}
	\item \textbf{Map:} any Map that has been unlocked via Research or previous completions.
	\item \textbf{Difficulty:} a Difficulty level for that Map (if unlocked).
	\item \textbf{Wave Range (optional):} within a Map, specific Wave ranges or starting points may be selected if previously reached.
\end{itemize}

Maps are presented as a \textbf{fixed list of unlocked Maps}, not as a rotating offer. 
Map and Difficulty are selected separately: a Map has a base identity (theme, layout, enemy families), while its Difficulties add scaling and modifications.

% TODO: Later define exact UI structure (e.g. Map grid + Difficulty panel, or combined card layout).

\subsection{Wave System}

\subsubsection*{Wave Structure}

\begin{itemize}
	\item Each Map has a \textbf{fixed Wave structure}:
	\begin{itemize}
		\item total number of Waves,
		\item base composition per Wave (enemy types, ranks, tags),
		\item temporal pacing (early/mid/late Waves).
	\end{itemize}
	\item Wave blueprints are \textbf{finely configurable}:
	\begin{itemize}
		\item relative amount of basic vs. Magic vs. Elite enemies,
		\item which Waves can spawn Bosses,
		\item where higher-density or special packs appear.
	\end{itemize}
	\item Design guideline for early MVP Maps:
	\begin{itemize}
		\item early Waves: mostly basic enemies, low pressure,
		\item mid Waves: increasing number of Elite enemies and more complex packs,
		\item late Waves: Bosses and high-pressure encounters.
	\end{itemize}
\end{itemize}

\subsubsection*{Bosses and Champions}

\begin{itemize}
	\item \textbf{Bosses:}
	\begin{itemize}
		\item are boosted versions of regular units with significantly higher stats,
		\item have more skills and more complex behaviour,
		\item typically appear in later Waves of a Map, but are not restricted to a single final Wave.
	\end{itemize}
	\item \textbf{Champions:}
	\begin{itemize}
		\item are \emph{legendary Bosses} with distinct move sets, potentially phase-based encounters,
		\item are intended as \textbf{endgame gates}, not standard MVP content,
		\item can be used to gate higher Difficulties or Tier 3 recipes and similar high-value unlocks.
	\end{itemize}
	\item Any Wave can, in principle, contain special enemies (Boss, Champion, Elite clusters), 
	but MVP Maps will follow a more predictable escalation pattern to keep learning curves manageable.
\end{itemize}

% RISK: Champions with unique movesets are content-expensive; should be late in the production pipeline.

\subsection{Difficulties and Empowerment}

\subsubsection*{Map Difficulties}

\begin{itemize}
	\item Each Map has multiple \textbf{Difficulty levels}.
	\item Higher Difficulties:
	\begin{itemize}
		\item increase enemy stats (HP, damage),
		\item may unlock additional enemy types or tags,
		\item increase rewards (loot quantity/quality, XP, research progress).
	\end{itemize}
	\item Difficulties are unlocked by:
	\begin{itemize}
		\item completing the previous Difficulty of the same Map,
		\item and/or unlocking via Research (World Tree, Heroes Tree, etc.).
	\end{itemize}
	\item Lower Difficulties remain \textbf{relevant for farming}:
	\begin{itemize}
		\item each Map has partially \emph{map-specific enemy types},
		\item Research goals may require killing specific enemy families/tags that are more common on lower tiers.
	\end{itemize}
\end{itemize}

\subsubsection*{Empower Maps}

\begin{itemize}
	\item \textbf{Empower Maps} adds optional modifiers on top of base Difficulty:
	\begin{itemize}
		\item more monsters,
		\item increased enemy stats,
		\item additional modifiers expressed as tags (e.g. +Projectiles, +Speed, Leech Aura, etc.).
	\end{itemize}
	\item Empowerment \textbf{costs meta-materials/meta-currency} (e.g. Remains, special tokens):
	\begin{itemize}
		\item this ties risk/reward tuning to the economy,
		\item prevents spamming high-reward setups without cost.
	\end{itemize}
	\item For MVP, Empowerment is considered an \emph{endgame feature}:
	\begin{itemize}
		\item can be included with lower priority,
		\item base Difficulties should function without Empower being required.
	\end{itemize}
\end{itemize}

% RISK: Empower + Difficulties + Map variants can explode combinatorially.
% TODO: Limit the number of stacked modifiers for MVP to keep tuning manageable.

\subsection{Live Events (Future)}

Map Live-Events are concepted as a future extension, not part of the MVP baseline.

Two main variants are envisioned:

\begin{itemize}
	\item \textbf{Variant A – Kill-based Quick-Time Choices (Vampire Survivors-like):}
	\begin{itemize}
		\item Events triggered based on kill thresholds or Wave milestones.
		\item Present the player with quick choices (e.g. temporary buffs, extra enemies for extra loot).
	\end{itemize}
	\item \textbf{Variant B – POE-style Interaction Objects:}
	\begin{itemize}
		\item In-Map interactables such as strongboxes, rituals, abysses.
		\item On click, additional enemy waves or mini-encounters spawn with improved loot rewards.
	\end{itemize}
\end{itemize}

For MVP these remain out of scope or very limited prototypes; core progression and difficulty should work without Live-Events.

% RISK: Live-Events add another layer of complexity on top of Waves and Empowerment.
%       Best introduced after core Map/Wave/Difficulty loop is stable.

\subsection{Rewards and Progression}

\subsubsection*{Per-Wave Rewards}

\begin{itemize}
	\item Regular enemies drop:
	\begin{itemize}
		\item Parts,
		\item Remains,
		\item currencies (primarily Gold).
	\end{itemize}
	\item Drops are handled by the Loot System; Maps define when and where drops are triggered.
	\item After each successful Wave:
	\begin{itemize}
		\item XP for surviving heroes is granted,
		\item the player may pick up all remaining loot or exit back to the Hideout,
		\item progress is saved up to the last successful Wave.
	\end{itemize}
\end{itemize}

\subsubsection*{Map Completion Rewards}

\begin{itemize}
	\item Completing all Waves of a Map:
	\begin{itemize}
		\item marks the Map as completed on that Difficulty,
		\item unlocks higher Difficulties for that Map (if available),
		\item yields an additional completion reward:
		\begin{itemize}
			\item meta-currency (e.g. for Research or Empowerment),
			\item progress in Research or Battlepass-like tracks,
			\item chance to unlock new recipes, heroes or Map variants (depending on design).
		\end{itemize}
	\end{itemize}
	\item Failure on a Map (all heroes dead) forfeits loot and XP of the current Wave, 
	but does not roll back global progression (research, unlocked Maps, etc.).
\end{itemize}

% RISK: If completion rewards are much better than per-wave rewards, players may feel forced to push beyond comfort.
% TODO: Balance per-wave vs. completion rewards so early exits still feel respectable for low-time players.

\subsection{MVP Scope for the Map System}

For the initial playable version (MVP), the Map System should deliver:

\begin{itemize}
	\item \textbf{Around 10 Maps}:
	\begin{itemize}
		\item spread across 3 themes,
		\item each Map with distinct layout/obstacle configuration,
		\item with fixed but varied Wave structures and enemy mixes.
	\end{itemize}
	\item \textbf{Multiple Difficulties} per Map:
	\begin{itemize}
		\item at least 2--3 tiers of Difficulty for core Maps,
		\item fully playable without Empowerment.
	\end{itemize}
	\item \textbf{Basic Boss implementation:}
	\begin{itemize}
		\item boosted enemies with more stats and skills.
		\item Champions and multi-phase bosses can be postponed or introduced as late MVP/endgame pieces.
	\end{itemize}
	\item \textbf{Empower Maps:}
	\begin{itemize}
		\item optional, low-priority endgame feature for MVP,
		\item only a small, curated set of modifiers initially.
	\end{itemize}
	\item \textbf{Live-Events:}
	\begin{itemize}
		\item considered \emph{Future}, not core MVP.
		\item optional to prototype a single simple event type for testing.
	\end{itemize}
\end{itemize}

% TODO: Cross-check final MVP Map count and Difficulty tiers with production capacity and balancing schedule.
